% --------------------------------------------
% CHAPTER 1: INTRODUCTION
% --------------------------------------------
\setlength{\parindent}{0.5in}

\chapter{Introduction}

The motivation: One of the stated goals of the DOE NEAMS program is to develop efficient multiscale fuel performance and structural materials modeling that accurately simulates temperature while leveraging experumental data to support industry and regulatory advancement \autocite{TechnicalAreas}, The US government is similarly motivated, with multiple executive orders (EOs) and budget documents having been released that call for increased development and work to be done in the energy sector \autocite{ordersDeployingAdvancedNuclear2025}. 

In the nuclear power industy, we look at recent political interest in advanced nuclear technologies that has created a large push for non-fossil fuel power. Motivations include the ongoing need for decreased emissions, the large scale need for power in an increasingly digital world, and recent technological advancements that have increased the demand for power-intensive datacenters for computational power. All of these things necessitate development of bigger and better reactor design technologies to support the advance of nuclear power.

The problem: Simulation methods are a demonstrated tool for research and development. The nuclear industry is interested in Accelerated Fuel Qualification (AFQ) for the accelerated research and development of nuclear materials, such as novel nuclear fuel types and geometries \autocite{ordersDeployingAdvancedNuclear2025,matthewsMechanisticMultiscaleUncertainty2026,choiAcceleratedFuelQualification2026}. Part of AFQ is multiscale modeling and simulation, which allows for more precise and focused experimental work with an associated simulation model to assist in qualification. These physics based simulations have proven a valuable tool for aiding in experimental work and giving predictions for materials properties that cannot be directly observed within the reactor environment. Thermal conductivity is an extrinsic material property, that depends on the material geometry and temperature as much as the chemical composition and crystal structure. We cannot directly observe fission product impact on thermal conductivity within the fuel element during reactor activation, but we can simulate it. 

The solution:Using a multiscale method combining nanoscale (ab inito) and mesoscale (Boltzmann) simualtion methods we can increase the accuracy of these simulations, allowing for better prediction of these and other difficult properties. Even further works utilize phase-field and cluster dynamics to incorporate neutronics, allowing for an even better picture of what is happening inside of the fuel element. The thermal conductivity is important for the heat release rate of the fuel, preventing too high of temperatures from accumulating within the fuel element, away from the coolant. Additionally, the fission products release rate and buildup rates have a large effect on fuel swelling, neutron trapping, and fuel-cladding interactions. We study the thermal interfacial resistance of nuclear fuels with fission products in order to predict some of these properties, and advance the development of a model that can accurately use fundamental physics-based methods to capture the effect of fission products in fuel material.


The connection: I demonstrate how my work using nanoscale methods is laying the foundation for future multiscale methods, and how important this kind of work is to the accelerated development of advanced reactors (alluding to further discussion in chapter 4)


Motivations! What is the problem being solved? Why is it a problem? General idea of giants in this field to give scope. Industry interest, government interest, why would someone care?

Document Outline: The rest of this document gives an overview of several years of work that led to this final product. It is organized to help give non-computational materials scientists a general understanding of the topics at hand (Chapter ~\ref{chap:background}).


\clearpage

% TENSE: Use present tense when describing general truths, current problems, or accepted knowledge, stating objectives or actions of the thesis. Use past tense when referring to past studies or actions.

%Open with the policy imperative (EO #3 → 400 GW demand), pivot to the qualification bottleneck (Matthews #1 → 20–25 year timeline), then introduce multiscale computation as the solution pathway (#2, #4, #5), landing on the gap your thesis fills: thermodynamic property prediction via DFT/DFTB+ feeding mesoscale methods within this exact framework.